\section{E. The remaining conditions and what the graded comparison sees}

\begin{claim}{E.1}{Proved}{The extra orbifold compatibility is arithmetic}
The outer action of $G_\mathbb Q$ satisfies the compatibility of B.m3,
understood first modulo inner automorphisms of the target $\widehat B$.
After its allowed cusp normalization it satisfies the equations
(7),(10),(15),(47),(51),(54) of \S47, with the sign and conjugator
corrections specified in B.m5--B.m6. These are formulations and consequences
of one compatibility, not six independent conjectural constraints.
\end{claim}
% anchor:E.1
\manuscriptanchor{\S 41; fols.315--316 verso; scans54--57; \S 47; fols.415--425; scans229--250 (cote140-3)}

\begin{proof}
Let $\mathcal X$ be the orbifold whose coarse curve is
$\mathbb P^1_\mathbb Q\setminus\{0\}$, with stabilizer orders two at
$1$ and three at $\infty$. The coordinate $T=1728/j$ identifies it with
the rigidification of $\mathcal M_{1,1}$ by its generic elliptic
involution; this can equally be seen from the weighted projective
presentation with weights two and three. Its geometric fundamental
group is $\widehat B=C_2\widehat * C_3$.

There are two maps over $\mathbb Q$ relevant here. The level-two cover
from the Legendre coordinate is
\[
 U\longrightarrow\mathcal X,\qquad
 T=\frac{27\lambda^2(1-\lambda)^2}
          {4(1-\lambda+\lambda^2)^3};
\]
it induces the inclusion $F\hookrightarrow B$ and the action used in
the dictionary. The open immersion
$\mathbb P^1_\mathbb Q\setminus\{0,1,\infty\}\hookrightarrow\mathcal X$
induces the surjection $\varphi$, because a meridian about $1$ acquires
order two and a meridian about $\infty$ acquires order three. Identify
the two tripod coordinates over $\mathbb Q$; this is the identification
used in the manuscript. Functoriality of the arithmetic fundamental
group for the immersion gives a homomorphism of arithmetic extensions.
Conjugating by a lift of a Galois element proves equivariance up to
geometric conjugation in $\widehat B$. This proves the outer assertion
without assuming a description of $G_\mathbb Q$ by equations.

For the explicit normalized equations, write the images of the two
free generators using B.m1. The first equation determines the possible
conjugator up to the procyclic centralizer of $e=sr$; write that conjugator as $e^t$. Replacing $u_0$ by
$\operatorname{Int}(y^n)u_0$ changes it to
$\varphi(y^n)e^t y^{-n}=e^{t-n}$, so the allowed choice $n=t$
removes it. The centralizer of $e$ is procyclic: reduce in $S_3$,
remove a power of $e$, and use the procyclic centralizer of
$e^2=y$ in $\widehat F$. The second is then exactly
B.m3. Checking instead the other free basis gives the order-three
formula (47); reducing these equalities modulo $F$ gives (51), and
their product gives (54). This is the generator-by-generator reduction
of B.m6. No assertion that (I),(II),(III) imply this compatibility is
used in the argument.
\end{proof}

\begin{claim}{E.2}{Proved}{The two-evaluation finite test through levels two and three}
There are twelve generating cycle residues in the characteristic
$S_3=\PSL_2(\mathbb F_2)$ envelope and twenty-four in the characteristic
$A_4=\PSL_2(\mathbb F_3)$ envelope. Every one admits a gauge satisfying
the normalized additional equation B.m3 after evaluation in the
respective projective congruence group. Thus the extra equation removes
none of these enumerated residues at these two levels. In particular every
actual profinite $\GTz$ element admits an allowed gauge with
$\varphi^\bullet(\xi)\in\{1,s\}$ in $S_3$; this finite consequence of
scan233 is therefore derived from (I),(II) and the automorphism condition.
\end{claim}
% anchor:E.2
\manuscriptanchor{\S 47; fols.417--420; scans233--240}

\begin{proof}
This calculation uses two homomorphisms from the \emph{same}
characteristic envelope to $Q_n$, not an ill-defined application of
$\varphi$ to an element of $H_n$. They are
\[
 j:(x,y)\mapsto(L_1,L_0),\qquad
 v:(x,y)\mapsto(S,SR).
\]
For $n=3$ the characteristic envelope has order $2304$, exponent six,
and derived group of order $256$. The analogous data at $n=2$ are in
B.f3; the map $j$ is trivial there and is still a well-defined map of
the envelope. The program constructs both maps by Cayley-graph
propagation and checks every relation. It enumerates (I),(II) and
generation exactly, with $\lambda$ modulo twelve and $f$ in the
computed derived group. For each solution put
$\nu=(\lambda-1)/2$, $\xi=y^\nu f$ and, for every $a$ modulo six,
$\xi_a=y^a\xi x^{-a}$. It then tests
\[
 j(\xi_a)=v(\xi_a)S v(\xi_a)^{-1}S^{-1}L_1^\nu.
 \tag{E.2.1}
\]
The period six includes every possible profinite gauge residue.
The recorded sets of successful gauges are nonempty for every one of
the twelve and twenty-four solutions, respectively. All coefficients,
all solutions, and all six defects per solution are retained in
\texttt{finite\_phi.json}; \texttt{finite\_phi.py} is the complete
enumeration. The assertion is exactly this necessary finite test; it
does not state that every such residue has an infinite lift.

For the last assertion, every actual profinite pair projects into this
list: the envelope is characteristic, the induced map is an automorphism,
and $f$ projects into its derived group. Its multiplier modulo twelve
determines $\nu$ modulo six. At level two $j$ and $L_1$ are trivial, so
(E.2.1) says that $v(\xi_a)$ centralizes $S$. That centralizer in $S_3$
is $\{1,S\}$. A successful residue $a\bmod6$ lifts to an allowed
$a\in\widehat{\mathbb Z}$. This is only a finite gauge condition;
it need not solve all the higher or full profinite equations.
\end{proof}

\begin{claim}{E.3}{Open}{Does the pentagon imply the additional orbifold compatibility?}
With the precise groups and normalizations of B.m3 and B.m5, the
remaining statement is
\[
 \GT\ \subseteq\ C^\varphi\ \subseteq\ \GTz.
 \tag{E.3.1}
\]
The second inclusion is definitional and the arithmetic image lies in
$C^\varphi$ by E.1. In this return we prove neither the first inclusion
nor its negation. The direct finite attack E.2 finds no violation at
target levels two and three. The next simultaneous-evaluation target
not enumerated for this equation is $Q_4\simeq S_4$, whose characteristic
envelope has the explicitly constructed order $442368$. Its outer
shadow is arithmetic by B.s4.1, but that outer shadow alone does not
retain the normalizing multiplier and conjugators needed in (E.2.1).
\end{claim}
% anchor:E.3
\manuscriptanchor{\S 41; fols.315--316 verso; scans54--57; \S 47; fols.415--425; scans229--250 (cote140-3)}

\begin{proof}[Attack and precise remaining step]
The proof of B.m2 derives every \S46 loop and symmetry equation from
(I),(II), with an inverse change of variables. It stops before applying
$\varphi$ to an unknown profinite word. Applying the congruence
homomorphism also does not supply that missing value: $j$ and $v$
are different maps. E.2 repairs this problem at two explicit finite
targets by retaining both evaluations. At level four the next required
calculation is to range over the normalized pairs in the characteristic
$S_4$ envelope, retaining their residue multiplier, and test (E.2.1)
for every gauge; the normalizer calculation of B.s4.1 does not by itself
perform that test. Even a violating finite pair would still require a
compatible pentagon lift before it could refute the first inclusion
in (E.3.1). Conversely finitely many passing targets would not prove
the inverse-limit inclusion.

We also compared the actual presentation with
\citet[Theorems 1.1--1.4 and relations (III$'$),(IV)]{NakamuraSchneps2000root}.
That source constructs an arithmetic-containing subgroup of $\GT$
using additional genus-one relations, including a Kummer cocycle for
the roots of two. It does not give an identified substitution proving
(E.3.1). We therefore do not rename its subgroup as $C^\varphi$ or
claim an equivalence of the additional relations without the missing
comparison. This is the exact literature-comparison obstruction, in
addition to the explicit finite calculation still required.
\end{proof}

\begin{claim}{E.4}{Proved}{Why the additional compatibility has no faithful graded form here}
The full condition E.1 cannot be faithfully linearized by replacing
both occurrences of the group by the single completed free Lie
algebra of Part B. The characteristic-zero pro-unipotent completion
of $B=C_2*C_3$ is trivial. Its pro-$\ell$ completion is $C_2$ for
$\ell=2$, $C_3$ for $\ell=3$, and trivial for every other prime.
Consequently the rational Lie algebra obtained from that target is
zero in every case. A relative completion retaining the $S_3$ quotient
would be additional data, not this absolute completion.
\end{claim}
% anchor:E.4
\manuscriptanchor{\S 47; fols.415--420; scans229--240}

\begin{proof}
An element of finite order in a characteristic-zero unipotent group
is the identity: taking its logarithm makes its order times that
logarithm zero. Any homomorphism from $C_2*C_3$ to such a group
therefore kills both generators. For a pro-$\ell$ target a generator
of order prime to $\ell$ is killed; the one surviving cyclic factor
at $\ell=2$ or $3$ gives precisely the claimed completion by its
universal property. These finite groups have zero rational
Malcev Lie algebra. Meanwhile $v:F\to S_3$ in E.2 is nontrivial,
so it cannot factor through a connected unipotent algebraic group
into the component group of a relative completion. One can restrict
to $\ker(v)$ and retain its larger free Lie algebra and comparison
maps, but this changes the input and the filtration. No such replacement
is silently made in $W^{\rm lac}$.
\end{proof}

\begin{claim}{E.5}{Proved}{The sorted inventory has one unresolved compatibility family}
Within the explicitly stated equations of \S\S38--47, the loop and
symmetry equations are (I),(II); the congruence-kernel and square-class
conditions follow for the whole actual profinite $\GTz$ by
C.d1 and C.a15--C.a16; the quadratic signs follow algebraically by
C.d2; and the literal uniqueness and multiplier claims have the
corrections proved in C.r2 and Part D. The additional equations still
not derived here from $\GT$ form the orbifold-compatibility family
in the following table. They are recorded separately in the condition
CSV because they retain different data, and their dependence is made
explicit rather than presented as independent discoveries.
\end{claim}
% anchor:E.5
\manuscriptanchor{\S 46; fols.390--413; scans182--227; \S 47; fols.415--427; scans229--253; inventory of \S \S 38--47}

\begin{longtable}{p{.13\textwidth}p{.37\textwidth}p{.13\textwidth}p{.26\textwidth}}
\toprule
Entry&Condition and location&Graded class&Decision and attack\\\midrule\endhead
E01&Outer $\varphi$ compatibility; \S41 scans55--57 and \S47 scans229--231, (5),(7).& (c)&Arithmetic by E.1; $\GT$, $\GTz$ undecided as in E.3.\\
E02&Normalized equation $\xi=[\varphi(\xi),s]x^\nu$; scans234,240, (15),(30).&(c)&Equivalent normalized version; complete simultaneous tests at levels2,3 in E.2.\\
E04&Cyclic order-three formula and its $S_3$ reduction; scans244--248, (47),(51).&(c)&Equivalent free-basis formulation after normalization; same unresolved global comparison.\\
E05&The pair of equations (57),(58) when $\mu\equiv1\pmod3$; scan250.&(c)&Their conjunction follows from E01; no claim that either single equation suffices.\\
\bottomrule
\end{longtable}
Here (a) would mean a proved derivation on the rational graded
pentagon Lie algebra, (b) would mean only a conditional use of the
2026 dimension equality, and (c) means there is no faithful graded
form in the specified absolute completion. All four remaining entries
are of class (c), by E.4. No entry is assigned to (b) merely to conceal
an unavailable completion map. The special-derivation identity itself
belongs to (a), by the direct pentagon derivation B.g8. The conditional
equality and its exact hypotheses are isolated in B.g9. The finite
signs restrict to $+1$ on $\lambda=1$ and have vacuous differentials,
as proved in C.a14.

The order-two finite consequence of scan233, formerly designated E03,
is excluded from this unresolved list: E.2 derives its existential-gauge
form for the whole $\GTz$. It remains in the complete condition inventory.
This distinction prevents a successful finite test from being reported
as an unresolved condition, or from being inflated into full compatibility.

For the auxiliary labels on scans247--248, B.m6 proves
$i\equiv\nu\pmod2$, equivalently
$i=(1-\epsilon_2(\mu))/2$, by applying the sign character to
$\varphi(\xi)=\beta s^i$. Thus this order-two label is already
determined algebraically once compatibility holds. The remaining
question whether $j'\in C_3$ is a function of the multiplier is not
an additional equation already proved on Galois. The explicit
complex-conjugation case gives $(i,a_3,j')=(1,1,-1)$, while the
identity gives $(0,0,0)$. The dependence of $j'$ is retained in the
finite-normalization problem E.3, without a guessed character.

\begin{claim}{E.6}{Proved}{No inference from the graded equality to profinite equality is used}
Even if the claimed 2026 upper bound is correct, it cannot on its own
decide (E.3.1), remove the congruence kernel of B.f6, distinguish the
four central lifts of C.a10, or classify full profinite torsion. It
also does not turn the finite Lazard counterexample of B.g15 into an
element of $\GTz\setminus\GT$.
\end{claim}
% anchor:E.6
\manuscriptanchor{\S 46; fols.390--413; scans182--227; \S 47; fols.415--427; scans229--253; inventory of \S \S 38--47}

\begin{proof}
Each proposed inference needs a comparison that is absent. E.4 shows
that the orbifold target disappears in the absolute Lie construction.
B.f6 supplies a noncongruence finite quotient missing from the matrix
completion. The four lifts differ by finite characters with zero
rational tangent space. A finite group quotient or a weight-truncated
Lie algebra does not produce a compatible inverse system of lifts.
Finally B.f8 gives an explicit endomorphism which is an automorphism
in every nilpotent and pro-$\ell$ realization but not in the full
profinite one. These are concrete obstructions to the respective
inferences, independent of the truth of the rational dimension theorem.
\end{proof}
